% cSpell:disable

\usetheme[
	progressbar=frametitle,
	titleformat section=smallcaps
]{metropolis}

\ifnotes%
	\setbeameroption{show notes on second screen}
\fi

\title{Data-X Talk} % chktex 13

\subtitle{Path ORAM:\ An Extremely Simple Oblivious RAM Protocol~\cite{DBLP:journals/corr/abs-1202-5150}}

\date{Built from \href{https://git.dbogatov.org/bu/PathORAM-Talk/commit/\version}{\emph{\version}} on \today}

\author{Dmytro Bogatov \\ dmytro@bu.edu}

\institute{Boston University}

\logo{
	\includegraphics[
		width=1cm,
		height=1cm,
		keepaspectratio
	]{static/logo.eps}\hspace{325pt}\vspace{-30pt}
}

\makeatletter
	\def\beamer@framenotesbegin{% at beginning of slide
		\usebeamercolor[fg]{normal text} 
			\gdef\beamer@noteitems{}%
			\gdef\beamer@notes{}% 
	}
\makeatother


\newcommand{\BigO}[1]{\mathcal{O}\left(#1\right)}
\newcommand{\BigOmega}[1]{\Omega\left(#1\right)}
\newcommand{\RAM}{\textbf{RAM}}

\DeclarePairedDelimiter\ceil{\lceil}{\rceil}
\DeclarePairedDelimiter\floor{\lfloor}{\rfloor}

\lstset{
	mathescape=true,
	numbers=left,
	tabsize=4,
	numbersep=10pt,
	breaklines=true,
	showtabs=false,
	morekeywords={for,end,do,while,if,then,else,Applies,when,return,true,false,Action},
	basicstyle=\normalfont,
	columns=fullflexible,
	frame=leftline,
	xleftmargin=15pt,
	escapechar=|
}

\newsavebox{\mybox}
